\scanpage{9}
\mspage{8}
Alors on trouve un $\otimes$-foncteur \inserted{pleinement} fidèle
\diagram{\begin{tikzcd}[ampersand replacement=\&]
\pi_1\text{-}Q\text{-modules admissibles de rang fini sur }Q \arrow[r] \& \text{motifs sur }S
\end{tikzcd}}
qui est décrit par un \inserted{épi}morphisme
\[
G=\varprojlim G_\alpha\unsure{\twoheadrightarrow|\longrightarrow}\pi_1(S,\xi)
\]
et se factorise donc en un \uline{épimorphisme} canonique
\[
\boxed{\varprojlim_{\alpha}G_\alpha/G_\alpha^0\longrightarrow\pi_1(S,\xi)}\tag*{(*)}
\]
Soit $\pi'_1$ le premier mb. Alors \deleted{les $\pi'_1$-} \inserted{(supposant $K=Q$} pour simplifier) les $\pi'_1$-$Q$-modules correspondent à la sous-catégorie de $M(S)$ formée des motifs \inserted{$M$} tels que la sous-\inserted{$\otimes$-}catégorie engendrée par $M$ soit de ``longueur finie''. Utilisant conjecture de Tate (ou de Hodge si $S$ de car. générique nulle) on \inserted{peut} \unsure{voir|voit} que ce sont les motifs \deleted{\unsure{de}} de niveau \inserted{et poids} zéro, \inserted{i.e. les motifs effectifs de poids zéro,} i.e. que $(*)$ \uline{est un isomorphisme}. Les \deleted{résultats}\inserted{réflexions} précédents peuvent se \deleted{\unsure{décrire}}\inserted{développer} aussi (utilisant la notion de motif sur $S$ à coeff. dans $K$) si on ne suppose pas $K=Q$.

\deleted{\uline{NB} La bijectivité de \unsure{ce signe}}
% letterhead
% printed upside down, lower right
INSTITUT\\
DES HAUTES ETUDES SCIENTIFIQUES\\
91 - BURES-SUR-YVETTE\\
928 48-53

\uline{NB} Soit $S'\longrightarrow{}S$ un revêtement étale \inserted{\unsure{surj}}\deleted{de} $S$, \inserted{$S'$ connexe,} et supposons que $F$ provienne d'
% archive sheet number: 8

